\documentclass[journal]{IEEEtran}
% Supplementary Material to the main paper (main.tex). Every number is a macro from numbers.tex, as in the paper.
\usepackage[utf8]{inputenc}
\usepackage[T1]{fontenc}
\usepackage{amsmath,amssymb}
\usepackage{graphicx}
\usepackage{booktabs}
\usepackage{xcolor}
\usepackage{cite}
\usepackage{url}
\usepackage[hidelinks,hypertexnames=false]{hyperref}
\graphicspath{{figs/}}
% Generated by src/paper/paper_numbers.py -- do not edit by hand.
\newcommand{\inwsDelivered}{19,090}  % data/clean/inws_cycles.csv node 1 rows
\newcommand{\inwsMissing}{16,230}  % results/scoping/inws_outages.csv ALL.lost_slots
\newcommand{\inwsLoss}{46.0\%}  % inws_outages: lost/(lost+delivered)
\newcommand{\inwsOutages}{1,619}  % inws_outages ALL.events
\newcommand{\inwsOneSlotEvents}{47.3\%}  % inws_outages 1-slot event_share
\newcommand{\inwsOneSlotLost}{4.7\%}  % inws_outages 1-slot lost_share
\newcommand{\inwsDayEvents}{14}  % inws_outages >=1 day events
\newcommand{\inwsDayEventShare}{0.9\%}  % inws_outages >=1 day event_share
\newcommand{\inwsDayLost}{50.7\%}  % outage_ci.csv InWaterSense L=24 share_obs
\newcommand{\inwsDayLostLo}{29.4\%}  % outage_ci.csv bootstrap 2.5%
\newcommand{\inwsDayLostHi}{64.6\%}  % outage_ci.csv bootstrap 97.5%
\newcommand{\inwsSixLost}{68.5\%}  % outage_ci.csv InWaterSense L=6 share_obs
\newcommand{\inwsSixLostLo}{54.9\%}  % outage_ci.csv
\newcommand{\inwsSixLostHi}{77.2\%}  % outage_ci.csv
\newcommand{\inwsRhoNextOut}{$+0.27$}  % outage_ci.csv spearman
\newcommand{\inwsRhoNextUp}{$-0.20$}  % outage_ci.csv spearman
\newcommand{\weijaRhoNextOut}{$+0.02$}  % outage_ci.csv spearman
\newcommand{\weijaPNextOut}{0.88}  % outage_ci.csv spearman
\newcommand{\weijaPNextUp}{0.73}  % outage_ci.csv spearman
\newcommand{\inwsMaxOutageDays}{13.3}  % inws_outages.log max slots x 611 s
\newcommand{\weijaStreams}{9}  % weija_outages.csv rows
\newcommand{\weijaCondLoss}{15.4\%}  % weija_outages.csv conductivity loss
\newcommand{\weijaCondOutages}{42}  % weija_outages.csv conductivity events
\newcommand{\weijaCondDayLost}{82.9\%}  % weija_outages.csv conductivity
\newcommand{\weijaCondDayLostLo}{0.0\%}  % outage_ci.csv Weija
\newcommand{\weijaCondDayLostHi}{93.5\%}  % outage_ci.csv Weija
\newcommand{\weijaCondDayN}{3}  % outage_ci.csv Weija n outages >= 1 d
\newcommand{\weijaDayLostMin}{71.0\%}  % weija_outages.csv min over streams
\newcommand{\weijaDayLostMax}{100.0\%}  % weija_outages.csv max over streams
\newcommand{\weijaMaxOutageDays}{36.7}  % weija_outages.csv max
\newcommand{\loedFrames}{5,118,817}  % loed_full.log
\newcommand{\loedDevices}{8,503}  % loed_full.log
\newcommand{\loedGateways}{9}  % loed_full.log
\newcommand{\loedDevicesTwoHundred}{1,780}  % loed_full.log
\newcommand{\loedWellHeard}{179}  % loed_full.log
\newcommand{\loedRadioLoss}{9.6\%}  % loed_full.log
\newcommand{\loedCondLoss}{0.25}  % loed_full.log
\newcommand{\loedBursts}{56,276}  % loed_full.log
\newcommand{\loedBurstOne}{84\%}  % loed_full.log
\newcommand{\loedBurstFive}{99\%}  % loed_full.log (length 1 + 2-5)
\newcommand{\loedCoLoss}{19.2\%}  % loed_full.log
\newcommand{\loedCoLossBase}{18.2\%}  % loed_full.log
\newcommand{\loedCoLossLo}{16.7\%}  % loed_full.log
\newcommand{\loedCoLossHi}{19.5\%}  % loed_full.log
\newcommand{\loedFullGateways}{5}  % loed_full_gateways.csv availability >= 0.999
\newcommand{\uvaAvail}{98.72\%}  % uva_loss.log
\newcommand{\uvaSysOutages}{12}  % uva_loss.log
\newcommand{\uvaSysMedianH}{1.67}  % uva_loss.log
\newcommand{\uvaSysMaxH}{82.2}  % uva_loss.log
\newcommand{\uvaSysDayShare}{83\%}  % uva_loss.log
\newcommand{\uvaMissing}{58,211}  % uva_loss.log
\newcommand{\uvaSysShare}{28.4\%}  % uva_loss.log
\newcommand{\uvaRadioShare}{71.6\%}  % uva_loss.log
\newcommand{\uvaGwShare}{0.0\%}  % uva_loss.log
\newcommand{\uvaLoss}{4.5\%}  % uva_loss.log
\newcommand{\uvaRuns}{36,519}  % uva_loss.log
\newcommand{\uvaRunsOne}{89\%}  % uva_loss.log
\newcommand{\uvaGwCDays}{29.9}  % uva_loss.log gateway C outage hours / 24
\newcommand{\uvaRainDry}{9.18\%}  % uva_loss.log
\newcommand{\uvaRainWet}{9.21\%}  % uva_loss.log
\newcommand{\uvaRainPairs}{18}  % uva_loss.log
\newcommand{\uvaRainP}{0.64}  % uva_loss.log
\newcommand{\uvaSensors}{10}  % uva_loss_sensors.csv
\newcommand{\uvaSensorFourRadio}{3.08\%}  % uva_loss_sensors.csv sensor04 radio/(rx+missing)
\newcommand{\aicBestInws}{discrete lognormal}  % outage model AIC table
\newcommand{\aicGeomGapInws}{3,677}  % AIC(geometric) - AIC(best)
\newcommand{\aicBestWeija}{discrete power law}  % outage model AIC table
\newcommand{\aicGeomGapWeija}{120}  % AIC(geometric) - AIC(best)
\newcommand{\aicBestUva}{two-geometric mixture}  % outage model AIC table
\newcommand{\aicGeomGapUva}{43,915}  % AIC(geometric) - AIC(best)
\newcommand{\arqEightMin}{2.59$\times$}  % loss_handling.csv LoRaWAN arq/unconf energy, min over records x SF x p
\newcommand{\arqEightMax}{7.35$\times$}  % same, max
\newcommand{\arqThreeMin}{1.51$\times$}  % loss_handling.csv LoRaWAN arq3/unconf energy, min
\newcommand{\arqThreeMax}{2.94$\times$}  % same, max
\newcommand{\arqGainMin}{$+5.0$}  % arq - unconf delivered share at p_radio = 0.096, min (pp)
\newcommand{\arqGainMax}{$+7.9$}  % same, max (pp)
\newcommand{\backoffMin}{1.20$\times$}  % NB-IoT unconf / unconf_backoff energy, min over records x search time x p
\newcommand{\backoffMax}{6.53$\times$}  % same, max
\newcommand{\bufNbMin}{$-5.5$}  % NB-IoT arq_buffer vs unconf_backoff energy change, min (percent)
\newcommand{\bufNbMax}{$+15.3$}  % same, max
\newcommand{\bufLoMin}{$-14.6$}  % LoRaWAN arq_buffer vs unconf energy change, min (percent)
\newcommand{\bufLoMax}{$+25.6$}  % same, max
\newcommand{\batchSaveMin}{1.33$\times$}  % fair baseline / batch1h energy, min over all configurations
\newcommand{\batchSaveMax}{3.33$\times$}  % same, max
\newcommand{\batchWithinMin}{$-3.3$}  % batch1h - unconf within-1h share, min (pp)
\newcommand{\batchWithinMax}{$+7.3$}  % same, max (pp)
\newcommand{\nConfigs}{20}  % loss_handling.csv record x radio x p_radio
\newcommand{\inwSevenUnconfE}{61.6}  % InWaterSense SF7 unconf J/day
\newcommand{\inwSevenUnconfLife}{1.80}  % InWaterSense SF7 unconf lifetime (y)
\newcommand{\inwSevenArqE}{447.1}  % InWaterSense SF7 arq J/day
\newcommand{\inwSevenArqLife}{0.46}  % InWaterSense SF7 arq lifetime (y)
\newcommand{\inwSevenBatchE}{18.9}  % InWaterSense SF7 batch1h J/day
\newcommand{\inwSevenBatchLife}{2.65}  % InWaterSense SF7 batch1h lifetime (y)
\newcommand{\inwRecordLossSim}{46.1\%}  % replayed InWaterSense link: unconf lost share (record grid)
\newcommand{\chainCount}{9}  % results/revision/chains/chain_<9 chains>.csv (D5 rerun)
\newcommand{\chainSites}{27}  % unique sites in the 9 chain files
\newcommand{\chainPairsMin}{53}  % ordered pairs per parameter, min
\newcommand{\chainPairsMax}{57}  % ordered pairs per parameter, max
\newcommand{\chainOutages}{300}  % median number of simulated gaps per pair x param x length
\newcommand{\chainRtWinT}{72\%}  % chains T 24 h: share of pairs where a neighbour method beats own data (real time)
\newcommand{\chainRtGainT}{$+6.5$}  % chains T 24 h: median relative real-time gain, pairs with positive baseline
\newcommand{\chainBfWinT}{93\%}  % chains T 24 h: share of pairs where the neighbour improves bridging
\newcommand{\chainBfGainT}{$+23.4$}  % chains T 24 h: median relative bridging gain, pairs with positive baseline
\newcommand{\chainDielBeatsPersistT}{84\%}  % chains T 24 h: share of pairs where own diel beats persistence
\newcommand{\chainRtWinDo}{30\%}  % chains DO 24 h: share of pairs where a neighbour method beats own data (real time)
\newcommand{\chainRtGainDo}{$-8.8$}  % chains DO 24 h: median relative real-time gain, pairs with positive baseline
\newcommand{\chainBfWinDo}{70\%}  % chains DO 24 h: share of pairs where the neighbour improves bridging
\newcommand{\chainBfGainDo}{$+10.3$}  % chains DO 24 h: median relative bridging gain, pairs with positive baseline
\newcommand{\chainDielBeatsPersistDo}{77\%}  % chains DO 24 h: share of pairs where own diel beats persistence
\newcommand{\chainExclSite}{14206694}  % the single site whose pairs are excluded
\newcommand{\chainRtWinPh}{9\%}  % chains pH 24 h: share of pairs where a neighbour method beats own data (real time)
\newcommand{\chainRtGainPh}{$-6.6$}  % chains pH 24 h: median relative real-time gain, pairs with positive baseline
\newcommand{\chainBfWinPh}{38\%}  % chains pH 24 h: share of pairs where the neighbour improves bridging
\newcommand{\chainBfGainPh}{$-0.2$}  % chains pH 24 h: median relative bridging gain, pairs with positive baseline
\newcommand{\chainDielBeatsPersistPh}{42\%}  % chains pH 24 h: share of pairs where own diel beats persistence
\newcommand{\chainRtWinEc}{23\%}  % chains EC 24 h: share of pairs where a neighbour method beats own data (real time)
\newcommand{\chainRtGainEc}{$-1.9$}  % chains EC 24 h: median relative real-time gain, pairs with positive baseline
\newcommand{\chainBfWinEc}{25\%}  % chains EC 24 h: share of pairs where the neighbour improves bridging
\newcommand{\chainBfGainEc}{$-2.2$}  % chains EC 24 h: median relative bridging gain, pairs with positive baseline
\newcommand{\chainDielBeatsPersistEc}{5\%}  % chains EC 24 h: share of pairs where own diel beats persistence
\newcommand{\chainGainExcluded}{3}  % 24 h pairs excluded from relative gains (zero baseline error)
\newcommand{\chainTolOneH}{97\%}  % chains: median over pairs x params of the best own-data real-time method within tolerance, 1 h
\newcommand{\chainTolDayMed}{31\%}  % chains: median over pairs x params of the best of ALL methods within tolerance, 24 h
\newcommand{\chainTolDayQ}{63\%}  % chains: 75th percentile of the same, 24 h
\newcommand{\eNbUnconf}{0.31}  % params.NBIOT N211 rai200 (J, module)
\newcommand{\eNbConf}{0.33}  % params.NBIOT N211 rai400 (J, module)
\newcommand{\eNbSearch}{2.22}  % P_RX_N211 x 10 s (J, module)
\newcommand{\pRxNb}{0.222}  % loss_handling.P_RX_N211 (W)
\newcommand{\eSfSevenUnconf}{0.244}  % model.lora_report_energy SF7 20 B unconfirmed (J, radio)
\newcommand{\eSfSevenConf}{0.164}  % model.lora_report_energy SF7 20 B confirmed success (J, radio)
\newcommand{\eSfTwelveUnconf}{0.729}  % model.lora_report_energy SF12 20 B unconfirmed (J, radio)
\newcommand{\eSfTwelveConf}{0.858}  % model.lora_report_energy SF12 20 B confirmed success (J, radio)
\newcommand{\rpfNb}{42}  % Radio(NB-IoT N211 RAI).readings_per_frame()
\newcommand{\rpfSfSeven}{18}  % Radio(LoRaWAN SF7).readings_per_frame()
\newcommand{\rpfSfTwelve}{3}  % Radio(LoRaWAN SF12).readings_per_frame()
\newcommand{\batteryKJ}{87.9}  % params.BATTERIES 6600 mAh x 3.7 V (kJ)
\newcommand{\hostOnMW}{69}  % params.P_HOST_ON (mW)
\newcommand{\weijaCondMaxDays}{7.5}  % weija_outages.csv conductivity max_outage_d
\newcommand{\weijaCondStepMin}{32.5}  % weija_outages.csv conductivity step_min
\newcommand{\inwsRhoP}{$p<0.001$}  % outage_ci.csv InWaterSense spearman, both p < 0.001 (asserted)
\newcommand{\weijaRhoNextUp}{$-0.06$}  % outage_ci.csv Weija spearman outage vs next up-run
\newcommand{\weijaCondSixLost}{88.9\%}  % outage_ci.csv Weija L=6 share_obs
\newcommand{\geomDayInws}{less than 0.001\%}  % outage_models.csv geom share24h_model InWaterSense (asserted < 1e-5)
\newcommand{\geomSixInws}{11.4\%}  % outage_models.csv geom share6h_model InWaterSense
\newcommand{\geomDayWeija}{36.9\%}  % outage_models.csv geom share24h_model Weija
\newcommand{\uvaSixLost}{27.1\%}  % uva_loss_runs.csv share of lost frames in runs >= 6 h
\newcommand{\uvaDayLost}{23.6\%}  % uva_loss_runs.csv share of lost frames in runs >= 24 h
\newcommand{\loedCondRatio}{2.6}  % loed_full.log P(loss|loss) / pooled radio loss
\newcommand{\aicWeijaLnGap}{1.7}  % outage_models.csv Weija AIC(lognormal) - AIC(power law)
\newcommand{\bufSevenMin}{$-14.6$}  % LoRaWAN SF7 arq_buffer vs unconf energy change, min over records x p (percent)
\newcommand{\bufSevenMax}{$+0.6$}  % same SF7, max (percent)
\newcommand{\bufTwelveMin}{$+7.8$}  % LoRaWAN SF12 arq_buffer vs unconf energy change, min over records x p (percent)
\newcommand{\bufTwelveMax}{$+25.6$}  % same SF12, max (percent)
\newcommand{\nbSearchTimes}{2.2, 10, 60}  % loss_handling.csv NB-IoT radio names (search time, s)
\newcommand{\inwSevenBufE}{52.7}  % InWaterSense SF7 arq_buffer J/day
\newcommand{\inwSevenBufLife}{1.93}  % InWaterSense SF7 arq_buffer lifetime (y)
\newcommand{\inwSevenArqThreeE}{175.5}  % InWaterSense SF7 arq3 J/day
\newcommand{\chainNaiveWinDo}{77\%}  % chains DO 24 h: share of pairs where a neighbour method beats persistence
\newcommand{\chainPersistDo}{9.07}  % chains DO 24 h: median over pairs of persistence error
\newcommand{\chainDielDo}{4.33}  % chains DO 24 h: median over pairs of own-diel error
\newcommand{\chainBestNbDo}{5.13}  % chains DO 24 h: median over pairs of best neighbour error
\newcommand{\chainNaiveWinPh}{36\%}  % chains pH 24 h: share of pairs where a neighbour method beats persistence
\newcommand{\chainPersistPh}{0.175}  % chains pH 24 h: median over pairs of persistence error
\newcommand{\chainDielPh}{0.125}  % chains pH 24 h: median over pairs of own-diel error
\newcommand{\chainBestNbPh}{0.149}  % chains pH 24 h: median over pairs of best neighbour error
\newcommand{\chainGaps}{1, 3, 6, 12, 24, 48, 72}  % chain files H_h values (h)
\newcommand{\weijaManyDecimals}{10.4\%}  % Weija conductivity, non-imputed: share of values with > 3 decimals
\newcommand{\weijaMaxDecimals}{14}  % Weija conductivity, non-imputed: most decimals in a value
\newcommand{\hostSleepMW}{0.12}  % params.P_HOST_SLEEP (mW)
\newcommand{\tSenseS}{5}  % params.T_SENSE typical (s awake per sample)
\newcommand{\nbPsmUW}{9.35}  % N211 PSM sleep power (uW)
\newcommand{\loraSleepUW}{148}  % LoRa radio sleep power (uW)
\newcommand{\vNode}{3.7}  % params.V_NODE (V)
\newcommand{\gofPmin}{0.005}  % bootstrap KS: smallest attainable p with B = 200
\newcommand{\gofGeomDInws}{0.48}  % gof_models InWaterSense geometric KS D
\newcommand{\gofBestDInws}{0.027}  % gof_models InWaterSense smallest KS D of the others
\newcommand{\gofGeomDUva}{0.26}  % gof_models UVA geometric KS D
\newcommand{\gofBestDUva}{0.008}  % gof_models UVA smallest KS D of the others
\newcommand{\gofWeijaLnP}{0.37}  % gof_models Weija EC lognormal p
\newcommand{\gofWeijaPlP}{0.64}  % gof_models Weija EC power-law p
\newcommand{\gofInwsN}{1,619}  % gof_models InWaterSense gaps
\newcommand{\gofUvaN}{36,519}  % gof_models UVA runs
\newcommand{\deepBridgeOneSaits}{65\%}  % deep: SAITS bridging 1 h, share within tolerance
\newcommand{\deepBridgeOneStem}{62\%}  % deep: StemGNN bridging 1 h, share within tolerance
\newcommand{\deepBridgeDaySaits}{25\%}  % deep: SAITS bridging 24 h, share within tolerance
\newcommand{\deepBridgeDayStem}{23\%}  % deep: StemGNN bridging 24 h, share within tolerance
\newcommand{\deepRtOneSaits}{57\%}  % deep: SAITS realtime 1 h, share within tolerance
\newcommand{\deepRtOneStem}{52\%}  % deep: StemGNN realtime 1 h, share within tolerance
\newcommand{\deepRtDaySaits}{18\%}  % deep: SAITS realtime 24 h, share within tolerance
\newcommand{\deepRtDayStem}{16\%}  % deep: StemGNN realtime 24 h, share within tolerance
\newcommand{\deepBridgeOneInterp}{97\%}  % deep: interpolation 1 h bridging, same gaps
\newcommand{\deepBridgeDayInterp}{31\%}  % deep: interpolation 24 h bridging, same gaps
\newcommand{\deepRtOneDiel}{91\%}  % deep: own diel 1 h real time, same gaps
\newcommand{\deepRtDayDiel}{29\%}  % deep: own diel 24 h real time, same gaps
\newcommand{\deepDayDoSaits}{4.3}  % deep: SAITS 24 h bridging DO median max error (%sat)
\newcommand{\deepDayDoInterp}{6.7}  % deep: interpolation, same
\newcommand{\deepDayPhSaits}{0.126}  % deep: SAITS 24 h bridging pH median max error
\newcommand{\deepDayPhInterp}{0.148}  % deep: interpolation, same
\newcommand{\deepDayBestAny}{55\%}  % deep file: best within-tolerance share of any estimator, any parameter, 24 h
\newcommand{\deepGaps}{48,455}  % deep: gap x parameter evaluations per model
\newcommand{\mcDraws}{100}  % energy_mc draws
\newcommand{\mcArqDearer}{100\%}  % energy_mc: LoRaWAN draws with ARQ dearer than unconfirmed
\newcommand{\mcArqLow}{1.8$\times$}  % energy_mc: LoRaWAN ARQ/unconf, lowest 2.5% quantile
\newcommand{\mcBatchCheapOther}{98\%}  % energy_mc: SF7/NB-IoT draws with batch cheaper than the fair baseline, min
\newcommand{\mcBatchCheapTwelveMin}{85\%}  % energy_mc: SF12 share, min over records
\newcommand{\mcBatchCheapTwelveMax}{93\%}  % energy_mc: SF12 share, max
\newcommand{\mcBatchTwelveHi}{1.24$\times$}  % energy_mc: SF12 batch/fair 97.5% quantile, max
\newcommand{\mcBufMedMin}{0.81$\times$}  % energy_mc: buffer/fair median, min over record x radio
\newcommand{\mcBufMedMax}{1.10$\times$}  % same, max
\newcommand{\mcBufLo}{0.50$\times$}  % energy_mc: buffer/fair 2.5% quantile, min
\newcommand{\mcBufHi}{2.08$\times$}  % energy_mc: buffer/fair 97.5% quantile, max
\newcommand{\mcNbArqMed}{1.01$\times$}  % energy_mc: NB-IoT ARQ/unconf median
\newcommand{\ugwLoraRatio}{3$\times$}  % Ugwuanyi 2021 Tab. 5 LoRaWAN 1.19 J vs model 0.448 J (44 B, SF7, host incl.)
\newcommand{\ugwNbRatio}{40$\times$}  % Ugwuanyi 2021 Tab. 5 NB-IoT 21.4 J vs model 0.539 J (44 B, RAI, host incl.)
\newcommand{\weijaEcMissing}{916}  % gap_cause_summary.json weija
\newcommand{\weijaNodeSilentShare}{84\%}  % gap_cause_summary.json weija
\newcommand{\weijaOtherStreamShare}{16\%}  % gap_cause_summary.json
\newcommand{\weijaLongAliveShare}{6\%}  % gap_cause_summary.json: slots of v0.1 gaps >= 6 h with another stream
\newcommand{\weijaNodeLoss}{13.9\%}  % gap_cause_weija_records.csv node record loss
\newcommand{\weijaNodeGaps}{31}  % gap_cause_weija_records.csv node record gaps
\newcommand{\weijaNodeDayLost}{81\%}  % gap_cause_weija_records.csv node record S(1 d)
\newcommand{\inwsWindowShare}{57\%}  % gap_cause_summary.json inws: lost slots inside node 2's window
\newcommand{\inwsNodeTwoSilent}{79\%}  % gap_cause_summary.json inws
\newcommand{\inwsSharedLong}{74\%}  % gap_cause_summary.json inws pi, gaps >= 6 h
\newcommand{\inwsSpecificLong}{26\%}  % 1 - pi, gaps >= 6 h
\newcommand{\inwsSpecificShort}{63\%}  % 1 - pi, gaps < 6 h
\newcommand{\offExpectedInws}{83\%}  % loss_B.csv InWaterSense SF7 f=0.5 buffer: delivered/expected
\newcommand{\offExpectedWeija}{92\%}  % loss_B.csv Weija (node record) SF7 f=0.5 buffer: delivered/expected
\newcommand{\offBufFairZero}{0.85}  % loss_B.csv InWaterSense SF7 buffer/unconf energy, f=0
\newcommand{\offBufFairHalf}{1.02}  % loss_B.csv InWaterSense SF7 buffer/unconf energy, f=0.5
\newcommand{\offRealisedWeija}{0.49}  % loss_B.csv: realised node-off share for nominal f = 0.1 (Weija has few long gaps)
\newcommand{\arqGainFullMin}{0.0}  % loss_A.csv InWaterSense LoRaWAN ARQ gain (pp), full, p_radio 0, SF7/SF12, all variants
\newcommand{\arqGainFullMax}{0.0}  % loss_A.csv same, max
\newcommand{\arqGainUniformMin}{0.1}  % loss_A.csv InWaterSense LoRaWAN ARQ gain (pp), uniform, p_radio 0, SF7/SF12, all variants
\newcommand{\arqGainUniformMax}{4.6}  % loss_A.csv same, max
\newcommand{\arqGainBurstMin}{6.5}  % loss_A.csv InWaterSense LoRaWAN ARQ gain (pp), burst, p_radio 0, SF7/SF12, all variants
\newcommand{\arqGainBurstMax}{7.7}  % loss_A.csv same, max
\newcommand{\arqLostPp}{46}  % loss_A.csv InWaterSense unconfirmed: lost share (pp)
\newcommand{\arqGainWeijaMax}{0.9}  % loss_A.csv Weija records: largest ARQ gain (pp) over models
\newcommand{\arqBufRatioMin}{0.67}  % loss_A.csv InWaterSense LoRaWAN buffer/unconf energy over models
\newcommand{\arqBufRatioMax}{1.08}  % same, max
\newcommand{\arqRatioMin}{2.3$\times$}  % loss_A.csv InWaterSense LoRaWAN ARQ/unconf energy over models, variants
\newcommand{\arqRatioMax}{7.3$\times$}  % same, max
\newcommand{\arqVOneZeroTwo}{7.3$\times$}  % loss_A.csv InWaterSense SF7 full: 8 tx, LoRaWAN 1.0.2 DR stepping
\newcommand{\arqVOneZeroFour}{6.8$\times$}  % loss_A.csv InWaterSense SF7 full: NbTrans 8, LoRaWAN 1.0.4
\newcommand{\ramLossTwo}{24\%}  % loss_C.csv InWaterSense RAM 2 KB: readings lost (max over radios/strategies)
\newcommand{\ramLossEight}{10\%}  % loss_C.csv InWaterSense RAM 8 KB: readings lost (max over radios/strategies)
\newcommand{\maxBufferKB}{22}  % loss_C.csv InWaterSense largest buffer (KB)
\newcommand{\flashJDayMax}{0.03}  % loss_C.csv: largest flash energy (J/day), worst case
\newcommand{\radioJDayMin}{9}  % loss_C.csv: smallest radio energy (J/day)
\newcommand{\drainTwelveMedH}{1.2}  % loss_C.csv InWaterSense SF12 buffer, duty-cycle drain: median clearance (h)
\newcommand{\drainTwelveMaxD}{63}  % same, longest clearance (days)
\newcommand{\drainTwelveLiveP}{25}  % same, p95 delay of live readings during backlogs (days)
\newcommand{\drainSevenMaxH}{64}  % loss_C.csv InWaterSense SF7 duty-cycle drain: longest clearance (h)
\newcommand{\kEverySlotSeven}{2.2$\times$}  % sweep_k_batch.csv InWaterSense SF7: probe every slot / K=6 h
\newcommand{\kEverySlotNb}{3.7$\times$}  % sweep_k_batch.csv InWaterSense NB-IoT: same
\newcommand{\kSixToDaySave}{12\%}  % sweep: SF7 energy saving K 6 -> 24 h
\newcommand{\kPninetySix}{20}  % sweep: SF7 p90 recovery latency at K = 6 h (h)
\newcommand{\kPninetyDay}{61}  % sweep: SF7 p90 recovery latency at K = 24 h (h)
\newcommand{\uvaMissThouVsFive}{2\%}  % lorawan_uva_counters.csv: missing frames, threshold 1,000 vs 5,000
\newcommand{\uvaDecreases}{1}  % lorawan_uva_counters.csv: frame-counter decreases
\newcommand{\uvaBinOutMax}{14}  % lorawan_uva_bins.csv: system outages at 5-min bins
\newcommand{\uvaBinOutMin}{8}  % lorawan_uva_bins.csv: system outages at 30-min bins
\newcommand{\uvaBinDayMin}{82\%}  % lorawan_uva_bins.csv: down-time share in >= 1 d outages, min over bins
\newcommand{\uvaBinDayMax}{84\%}  % same, max
\newcommand{\uvaBinSysMin}{28.0\%}  % lorawan_uva_bins.csv: system share of missing frames, min
\newcommand{\uvaBinSysMax}{28.6\%}  % same, max
\newcommand{\loedLossHundred}{10.0\%}  % lorawan_loed_counters.csv well-heard devices, reset threshold 100
\newcommand{\loedLossThousand}{10.9\%}  % same, threshold 1,000 (the analysis)
\newcommand{\loedResetDevices}{81}  % well-heard devices with a frame-counter reset
\newcommand{\loedWellHeardN}{179}  % well-heard devices
\newcommand{\usgsExcT}{0.20}  % tolerance_summary.csv T: USGS excellent (median over sites)
\newcommand{\tolTwoT}{0.20}  % tolerance_summary.csv T: L2 (median over sites)
\newcommand{\dayRangeT}{1.10}  % tolerance_summary.csv T: median daily range
\newcommand{\ampFiveT}{0.74}  % tolerance_summary.csv T: median 5 sigma
\newcommand{\ampEightT}{1.19}  % tolerance_summary.csv T: median 8 sigma
\newcommand{\realFiveT}{100\%}  % tolerance_summary.csv T: sites with a real 3-h change >= 5 sigma
\newcommand{\realEightT}{83\%}  % tolerance_summary.csv T: sites with a real 3-h change >= 8 sigma
\newcommand{\usgsExcDO}{4.7}  % tolerance_summary.csv DO: USGS excellent (median over sites)
\newcommand{\tolTwoDO}{2.0}  % tolerance_summary.csv DO: L2 (median over sites)
\newcommand{\dayRangeDO}{10.0}  % tolerance_summary.csv DO: median daily range
\newcommand{\ampFiveDO}{7.4}  % tolerance_summary.csv DO: median 5 sigma
\newcommand{\ampEightDO}{11.9}  % tolerance_summary.csv DO: median 8 sigma
\newcommand{\realFiveDO}{100\%}  % tolerance_summary.csv DO: sites with a real 3-h change >= 5 sigma
\newcommand{\realEightDO}{97\%}  % tolerance_summary.csv DO: sites with a real 3-h change >= 8 sigma
\newcommand{\usgsExcEC}{6.3}  % tolerance_summary.csv EC: USGS excellent (median over sites)
\newcommand{\tolTwoEC}{4.2}  % tolerance_summary.csv EC: L2 (median over sites)
\newcommand{\dayRangeEC}{12.0}  % tolerance_summary.csv EC: median daily range
\newcommand{\ampFiveEC}{7.4}  % tolerance_summary.csv EC: median 5 sigma
\newcommand{\ampEightEC}{11.9}  % tolerance_summary.csv EC: median 8 sigma
\newcommand{\realFiveEC}{93\%}  % tolerance_summary.csv EC: sites with a real 3-h change >= 5 sigma
\newcommand{\realEightEC}{93\%}  % tolerance_summary.csv EC: sites with a real 3-h change >= 8 sigma
\newcommand{\usgsExcPh}{0.20}  % tolerance_summary.csv pH: USGS excellent (median over sites)
\newcommand{\tolTwoPh}{0.10}  % tolerance_summary.csv pH: L2 (median over sites)
\newcommand{\dayRangePh}{0.25}  % tolerance_summary.csv pH: median daily range
\newcommand{\ampFivePh}{0.50}  % tolerance_summary.csv pH: median 5 sigma
\newcommand{\ampEightPh}{0.80}  % tolerance_summary.csv pH: median 8 sigma
\newcommand{\realFivePh}{60\%}  % tolerance_summary.csv pH: sites with a real 3-h change >= 5 sigma
\newcommand{\realEightPh}{37\%}  % tolerance_summary.csv pH: sites with a real 3-h change >= 8 sigma
\newcommand{\netZeroJitterN}{50}  % netsim.csv: smallest cell (SF7) in which, without jitter, under 5% of nodes recover within 72 h
\newcommand{\netSevenMedMin}{2}  % netsim SF7, 1-h jitter: median node clearance, min (h)
\newcommand{\netSevenMedMax}{9}  % same, max
\newcommand{\netSevenAllMax}{42}  % netsim SF7: last node clears (h), max over settings
\newcommand{\netSevenEnergyMax}{12.5$\times$}  % netsim SF7, 500 nodes: energy per node vs alone, max
\newcommand{\netTwelveNoAck}{29\%}  % netsim SF12: received but unacknowledged, max
\newcommand{\netTwelveRxOne}{1.0\%}  % netsim SF12: gateway RX1 duty use, max (limit 1%)
\newcommand{\netTwelveCleared}{93\%}  % netsim SF12, 1-d outage, 100 nodes: share recovered in 72 h
\newcommand{\recDayAll}{31\%}  % reconstruct_bootstrap.csv: best-of-all within tolerance, 24 h, pooled
\newcommand{\recDayAllLo}{18\%}  % chain-blocked 95% CI low
\newcommand{\recDayAllHi}{49\%}  % CI high
\newcommand{\recDayT}{9\%}  % reconstruct_bootstrap.csv T 24 h best-of-all
\newcommand{\recDayPh}{49\%}  % reconstruct_bootstrap.csv pH 24 h best-of-all
\newcommand{\recDayDO}{22\%}  % reconstruct_bootstrap.csv DO 24 h best-of-all
\newcommand{\recDayEC}{56\%}  % reconstruct_bootstrap.csv EC 24 h best-of-all
\newcommand{\lagPosDown}{64\%}  % pairs_spatial DO: downstream neighbours with lag > 0
\newcommand{\lagPosUp}{68\%}  % pairs_spatial DO: upstream neighbours with lag > 0
\newcommand{\pairKmMin}{2.8}  % pairs_spatial distance min (km)
\newcommand{\pairKmMax}{300}  % max (km)
\newcommand{\pairKmMed}{42}  % median (km)
\newcommand{\chOneDev}{12/12}  % an2023_verdicts.json CH1
\newcommand{\chTwoDev}{12/12}  % an2023_verdicts.json CH2
\newcommand{\chFourDev}{38/48}  % an2023_verdicts.json CH4
\newcommand{\chFiveDev}{48/48}  % an2023_verdicts.json CH5
\newcommand{\chSixDev}{12/12}  % an2023_verdicts.json CH6
\newcommand{\chSevenDev}{12/12}  % an2023_verdicts.json CH7
\newcommand{\chNineLOneDev}{48/48}  % an2023_verdicts.json CH9_L1
\newcommand{\chNineLThreeDev}{48/48}  % an2023_verdicts.json CH9_L3
\newcommand{\chTenLOneDev}{12/12}  % an2023_verdicts.json CH10_L1
\newcommand{\chTenLThreeDev}{12/12}  % an2023_verdicts.json CH10_L3
\newcommand{\plausibleDev}{99.8\%}  % an2023_verdicts.json: injected events within plausible ranges
\newcommand{\designSites}{30}  % design_site_eligibility.csv status ok
\newcommand{\designCandidates}{39}  % design_site_eligibility.csv rows
\newcommand{\hFourFailMinDev}{$-10.8$}  % an2023_picture L2 worst failing event - arq (pp)
\newcommand{\hFourFailMaxDev}{$-5.1$}  % an2023_picture L2 mildest failing event - arq (pp)
\newcommand{\hFiveDiffMinDev}{$-0.02$}  % an2023_picture L2 bounded - arq, min (pp)
\newcommand{\hFiveDiffMaxDev}{$+0.59$}  % an2023_picture L2 bounded - arq, max (pp)
\newcommand{\eventOverBatchMinDev}{1.00$\times$}  % an2023_energy event / batch, min
\newcommand{\eventOverBatchMaxDev}{1.02$\times$}  % an2023_energy event / batch, max
\newcommand{\arqOverEventMinDev}{1.41$\times$}  % an2023_energy arq / event, min
\newcommand{\arqOverEventMaxDev}{3.86$\times$}  % an2023_energy arq / event, max
\newcommand{\boundedOverBatchMinDev}{1.05$\times$}  % an2023_energy bounded / batch, min
\newcommand{\boundedOverBatchMaxDev}{1.34$\times$}  % an2023_energy bounded / batch, max
\newcommand{\arqOverBoundedMinDev}{1.32$\times$}  % an2023_energy arq / bounded, min
\newcommand{\arqOverBoundedMaxDev}{2.51$\times$}  % an2023_energy arq / bounded, max
\newcommand{\arqOverBoundedCiLowDev}{1.29$\times$}  % an2023_energy: lowest lower 95% CI bound of arq / bounded over link x radio
\newcommand{\levOneArqOverMinDev}{1.20$\times$}  % an2023_energy_levels L1 arq / bounded, min
\newcommand{\levOneArqOverMaxDev}{1.61$\times$}  % an2023_energy_levels L1 arq / bounded, max
\newcommand{\levTwoArqOverMinDev}{1.33$\times$}  % an2023_energy_levels L2 arq / bounded, min
\newcommand{\levTwoArqOverMaxDev}{2.54$\times$}  % an2023_energy_levels L2 arq / bounded, max
\newcommand{\levThreeArqOverMinDev}{1.40$\times$}  % an2023_energy_levels L3 arq / bounded, min
\newcommand{\levThreeArqOverMaxDev}{3.38$\times$}  % an2023_energy_levels L3 arq / bounded, max
\newcommand{\batchLagMinDev}{0.50}  % an2023_latency: batch - event (h), min over link x radio
\newcommand{\batchLagMaxDev}{0.50}  % an2023_latency: batch - event (h), max
\newcommand{\knownOneUnconfNoneDev}{33\%}  % an2023_known_within none SF7 unconf 1 h
\newcommand{\knownSixUnconfNoneDev}{57\%}  % an2023_known_within none SF7 unconf 6 h
\newcommand{\knownDayUnconfNoneDev}{57\%}  % an2023_known_within none SF7 unconf 24 h
\newcommand{\knownOneArqNoneDev}{33\%}  % an2023_known_within none SF7 arq_buffer 1 h
\newcommand{\knownSixArqNoneDev}{57\%}  % an2023_known_within none SF7 arq_buffer 6 h
\newcommand{\knownDayArqNoneDev}{57\%}  % an2023_known_within none SF7 arq_buffer 24 h
\newcommand{\knownOneBatchNoneDev}{21\%}  % an2023_known_within none SF7 batch1h 1 h
\newcommand{\knownSixBatchNoneDev}{57\%}  % an2023_known_within none SF7 batch1h 6 h
\newcommand{\knownDayBatchNoneDev}{57\%}  % an2023_known_within none SF7 batch1h 24 h
\newcommand{\knownOneEventNoneDev}{33\%}  % an2023_known_within none SF7 design 1 h
\newcommand{\knownSixEventNoneDev}{57\%}  % an2023_known_within none SF7 design 6 h
\newcommand{\knownDayEventNoneDev}{57\%}  % an2023_known_within none SF7 design 24 h
\newcommand{\knownOneBoundedNoneDev}{33\%}  % an2023_known_within none SF7 design_bounded_nf 1 h
\newcommand{\knownSixBoundedNoneDev}{57\%}  % an2023_known_within none SF7 design_bounded_nf 6 h
\newcommand{\knownDayBoundedNoneDev}{57\%}  % an2023_known_within none SF7 design_bounded_nf 24 h
\newcommand{\knownOneUnconfInwsDev}{12\%}  % an2023_known_within InWaterSense SF7 unconf 1 h
\newcommand{\knownSixUnconfInwsDev}{21\%}  % an2023_known_within InWaterSense SF7 unconf 6 h
\newcommand{\knownDayUnconfInwsDev}{21\%}  % an2023_known_within InWaterSense SF7 unconf 24 h
\newcommand{\knownOneArqInwsDev}{17\%}  % an2023_known_within InWaterSense SF7 arq_buffer 1 h
\newcommand{\knownSixArqInwsDev}{35\%}  % an2023_known_within InWaterSense SF7 arq_buffer 6 h
\newcommand{\knownDayArqInwsDev}{43\%}  % an2023_known_within InWaterSense SF7 arq_buffer 24 h
\newcommand{\knownOneBatchInwsDev}{11\%}  % an2023_known_within InWaterSense SF7 batch1h 1 h
\newcommand{\knownSixBatchInwsDev}{35\%}  % an2023_known_within InWaterSense SF7 batch1h 6 h
\newcommand{\knownDayBatchInwsDev}{44\%}  % an2023_known_within InWaterSense SF7 batch1h 24 h
\newcommand{\knownOneEventInwsDev}{18\%}  % an2023_known_within InWaterSense SF7 design 1 h
\newcommand{\knownSixEventInwsDev}{38\%}  % an2023_known_within InWaterSense SF7 design 6 h
\newcommand{\knownDayEventInwsDev}{44\%}  % an2023_known_within InWaterSense SF7 design 24 h
\newcommand{\knownOneBoundedInwsDev}{19\%}  % an2023_known_within InWaterSense SF7 design_bounded_nf 1 h
\newcommand{\knownSixBoundedInwsDev}{38\%}  % an2023_known_within InWaterSense SF7 design_bounded_nf 6 h
\newcommand{\knownDayBoundedInwsDev}{44\%}  % an2023_known_within InWaterSense SF7 design_bounded_nf 24 h
\newcommand{\chOneConf}{12/12}  % an2024_verdicts.json CH1
\newcommand{\chTwoConf}{12/12}  % an2024_verdicts.json CH2
\newcommand{\chFourConf}{39/48}  % an2024_verdicts.json CH4
\newcommand{\chFiveConf}{48/48}  % an2024_verdicts.json CH5
\newcommand{\chSixConf}{12/12}  % an2024_verdicts.json CH6
\newcommand{\chSevenConf}{12/12}  % an2024_verdicts.json CH7
\newcommand{\chNineLOneConf}{48/48}  % an2024_verdicts.json CH9_L1
\newcommand{\chNineLThreeConf}{48/48}  % an2024_verdicts.json CH9_L3
\newcommand{\chTenLOneConf}{12/12}  % an2024_verdicts.json CH10_L1
\newcommand{\chTenLThreeConf}{12/12}  % an2024_verdicts.json CH10_L3
\newcommand{\plausibleConf}{99.7\%}  % an2024_verdicts.json: injected events within plausible ranges
\newcommand{\confSites}{33}  % eligibility_2024.csv status ok
\newcommand{\hFourFailMinConf}{$-9.0$}  % an2024_picture L2 worst failing event - arq (pp)
\newcommand{\hFourFailMaxConf}{$-6.9$}  % an2024_picture L2 mildest failing event - arq (pp)
\newcommand{\hFiveDiffMinConf}{$-0.01$}  % an2024_picture L2 bounded - arq, min (pp)
\newcommand{\hFiveDiffMaxConf}{$+0.56$}  % an2024_picture L2 bounded - arq, max (pp)
\newcommand{\eventOverBatchMinConf}{1.00$\times$}  % an2024_energy event / batch, min
\newcommand{\eventOverBatchMaxConf}{1.02$\times$}  % an2024_energy event / batch, max
\newcommand{\arqOverEventMinConf}{1.41$\times$}  % an2024_energy arq / event, min
\newcommand{\arqOverEventMaxConf}{3.86$\times$}  % an2024_energy arq / event, max
\newcommand{\boundedOverBatchMinConf}{1.05$\times$}  % an2024_energy bounded / batch, min
\newcommand{\boundedOverBatchMaxConf}{1.28$\times$}  % an2024_energy bounded / batch, max
\newcommand{\arqOverBoundedMinConf}{1.34$\times$}  % an2024_energy arq / bounded, min
\newcommand{\arqOverBoundedMaxConf}{2.55$\times$}  % an2024_energy arq / bounded, max
\newcommand{\arqOverBoundedCiLowConf}{1.28$\times$}  % an2024_energy: lowest lower 95% CI bound of arq / bounded over link x radio
\newcommand{\levOneArqOverMinConf}{1.20$\times$}  % an2024_energy_levels L1 arq / bounded, min
\newcommand{\levOneArqOverMaxConf}{1.57$\times$}  % an2024_energy_levels L1 arq / bounded, max
\newcommand{\levTwoArqOverMinConf}{1.35$\times$}  % an2024_energy_levels L2 arq / bounded, min
\newcommand{\levTwoArqOverMaxConf}{2.58$\times$}  % an2024_energy_levels L2 arq / bounded, max
\newcommand{\levThreeArqOverMinConf}{1.40$\times$}  % an2024_energy_levels L3 arq / bounded, min
\newcommand{\levThreeArqOverMaxConf}{3.47$\times$}  % an2024_energy_levels L3 arq / bounded, max
\newcommand{\batchLagMinConf}{0.43}  % an2024_latency: batch - event (h), min over link x radio
\newcommand{\batchLagMaxConf}{0.50}  % an2024_latency: batch - event (h), max
\newcommand{\knownOneUnconfNoneConf}{34\%}  % an2024_known_within none SF7 unconf 1 h
\newcommand{\knownSixUnconfNoneConf}{58\%}  % an2024_known_within none SF7 unconf 6 h
\newcommand{\knownDayUnconfNoneConf}{58\%}  % an2024_known_within none SF7 unconf 24 h
\newcommand{\knownOneArqNoneConf}{34\%}  % an2024_known_within none SF7 arq_buffer 1 h
\newcommand{\knownSixArqNoneConf}{58\%}  % an2024_known_within none SF7 arq_buffer 6 h
\newcommand{\knownDayArqNoneConf}{58\%}  % an2024_known_within none SF7 arq_buffer 24 h
\newcommand{\knownOneBatchNoneConf}{21\%}  % an2024_known_within none SF7 batch1h 1 h
\newcommand{\knownSixBatchNoneConf}{58\%}  % an2024_known_within none SF7 batch1h 6 h
\newcommand{\knownDayBatchNoneConf}{58\%}  % an2024_known_within none SF7 batch1h 24 h
\newcommand{\knownOneEventNoneConf}{34\%}  % an2024_known_within none SF7 design 1 h
\newcommand{\knownSixEventNoneConf}{58\%}  % an2024_known_within none SF7 design 6 h
\newcommand{\knownDayEventNoneConf}{58\%}  % an2024_known_within none SF7 design 24 h
\newcommand{\knownOneBoundedNoneConf}{34\%}  % an2024_known_within none SF7 design_bounded_nf 1 h
\newcommand{\knownSixBoundedNoneConf}{58\%}  % an2024_known_within none SF7 design_bounded_nf 6 h
\newcommand{\knownDayBoundedNoneConf}{58\%}  % an2024_known_within none SF7 design_bounded_nf 24 h
\newcommand{\knownOneUnconfInwsConf}{12\%}  % an2024_known_within InWaterSense SF7 unconf 1 h
\newcommand{\knownSixUnconfInwsConf}{22\%}  % an2024_known_within InWaterSense SF7 unconf 6 h
\newcommand{\knownDayUnconfInwsConf}{22\%}  % an2024_known_within InWaterSense SF7 unconf 24 h
\newcommand{\knownOneArqInwsConf}{16\%}  % an2024_known_within InWaterSense SF7 arq_buffer 1 h
\newcommand{\knownSixArqInwsConf}{34\%}  % an2024_known_within InWaterSense SF7 arq_buffer 6 h
\newcommand{\knownDayArqInwsConf}{42\%}  % an2024_known_within InWaterSense SF7 arq_buffer 24 h
\newcommand{\knownOneBatchInwsConf}{11\%}  % an2024_known_within InWaterSense SF7 batch1h 1 h
\newcommand{\knownSixBatchInwsConf}{34\%}  % an2024_known_within InWaterSense SF7 batch1h 6 h
\newcommand{\knownDayBatchInwsConf}{43\%}  % an2024_known_within InWaterSense SF7 batch1h 24 h
\newcommand{\knownOneEventInwsConf}{18\%}  % an2024_known_within InWaterSense SF7 design 1 h
\newcommand{\knownSixEventInwsConf}{36\%}  % an2024_known_within InWaterSense SF7 design 6 h
\newcommand{\knownDayEventInwsConf}{44\%}  % an2024_known_within InWaterSense SF7 design 24 h
\newcommand{\knownOneBoundedInwsConf}{19\%}  % an2024_known_within InWaterSense SF7 design_bounded_nf 1 h
\newcommand{\knownSixBoundedInwsConf}{36\%}  % an2024_known_within InWaterSense SF7 design_bounded_nf 6 h
\newcommand{\knownDayBoundedInwsConf}{44\%}  % an2024_known_within InWaterSense SF7 design_bounded_nf 24 h
\newcommand{\designEventsConf}{7,548}  % 2024 injected events (all settings)
\newcommand{\detectStepFive}{67\%}  % 2024: median over sites of the share of step5 events detected
\newcommand{\detectStepEight}{91\%}  % 2024: median over sites of the share of step8 events detected
\newcommand{\detectRampFive}{13\%}  % 2024: median over sites of the share of ramp5 events detected
\newcommand{\detectRampEight}{65\%}  % 2024: median over sites of the share of ramp8 events detected
\newcommand{\detectAll}{58.2\%}  % 2024: share of injected events detected (pooled)
\newcommand{\faMedian}{0.8}  % 2024: median over sites of false alarms per parameter per 30 days
\newcommand{\faMean}{1.1}  % 2024: mean
\newcommand{\alarmsPerDay}{0.13}  % 2024: median node alarms per day (clean series)
\newcommand{\alarmsPerDayMean}{0.15}  % 2024: mean
\newcommand{\onsetUnconfKnownMax}{7.0\%}  % 2024 onset in outage: unconf known ever, max
\newcommand{\onsetBufKnownMin}{98.8\%}  % 2024: buffering known ever, min
\newcommand{\picInwsDoUnconf}{75.8\%}  % an2024 L2 InWaterSense SF7 DO unconf
\newcommand{\picInwsDoArq}{72.4\%}  % an2024 L2 InWaterSense SF7 DO arq_buffer
\newcommand{\nfRatioMin}{0.987}  % 2024 clean L2: newest-first / FIFO energy, min
\newcommand{\nfRatioMax}{1.007}  % 2024 clean L2: newest-first / FIFO energy, max
\newcommand{\sitesBothYears}{29}  % sites eligible in both years
\newcommand{\sitesOnlyConf}{4}  % sites eligible only in 2024
\newcommand{\hEightPreSite}{01467200}  % v0.1 pre_D4: the only site where newest-first fell short
\newcommand{\hEightPreMin}{99.998\%}  % v0.1 pre_D4: min within-tolerance share of newest-first, no outage
\newcommand{\fifoMinSfTwelve}{98.72\%}  % v0.1 summary_H8 FIFO min within-tolerance, SF12, no outage
\newcommand{\hFourPassVOne}{29/48}  % v0.1 summary_H4 (before the D5 fix)
\newcommand{\ackCostMin}{3\%}  % ack_variant: application / network energy - 1, min over link x radio
\newcommand{\ackCostMax}{18\%}  % same, max
\newcommand{\ackArqOverMin}{1.29$\times$}  % ack_variant: ARQ+buffer / bounded (application ACK), min
\newcommand{\ackArqOverMax}{2.00$\times$}  % same, max
\newcommand{\ackPicWorst}{$-0.02$}  % ack_variant: bounded (application ACK) - ARQ+buffer picture, worst median site (pp)

\makeatletter\newcommand{\tabin}[1]{\@@input #1 }\makeatother
\renewcommand{\thesection}{S\arabic{section}}
\renewcommand{\thetable}{S\arabic{table}}
\renewcommand{\thefigure}{S\arabic{figure}}
\renewcommand{\theequation}{S\arabic{equation}}

\begin{document}

\title{Supplementary Material for\\ ``Outage-Aware Reporting and Loss Recovery for Energy-Efficient River Water-Quality Telemetry over LPWANs''}
\author{Kofi Sarpong Adu-Manu and Stephen Kwaku Azongo}
\maketitle

This supplement gives the details behind the main paper: the delivery records and the cause of their gaps (Sec.~\ref{sec:s-records}), the gap-length models and their tests (Sec.~\ref{sec:s-models}), the energy model and its assumptions (Sec.~\ref{sec:s-energy}), further loss-handling results (Sec.~\ref{sec:s-handling}), the reconstruction method (Sec.~\ref{sec:s-reconstruct}), and the protocol and further results of the design evaluation (Sec.~\ref{sec:s-design}). References to the main paper are marked ``main paper''.

\section{Delivery records and the cause of gaps}
\label{sec:s-records}

\begin{figure*}[t]
\centering
\includegraphics[width=\textwidth]{fig_records}
\caption{Delivery timelines of the three records with long gaps: coloured spans are slots in which nothing was delivered (UVA: system outages only; its short radio losses are not shown). Each record is drawn on its own time axis, from its first slot.}
\label{fig:s-records}
\end{figure*}

Fig.~\ref{fig:s-records} shows when each record was down: many short gaps, and a few gaps lasting days.

\emph{InWaterSense.} The node sends each parameter as a separate message, seconds apart; the messages of one sampling instant are grouped into a reading, and a slot counts as delivered if a reading arrived for it.

\emph{Weija.} Only records not marked as imputed count as delivered. Even among them, \weijaManyDecimals{} of the conductivity values carry more than the three decimals the logger otherwise records (up to \weijaMaxDecimals), a sign of unflagged interpolation. Such values hide some short gaps, so counts of short gaps are lower bounds; long gaps are not affected.

\emph{UVA.} Missing frames are found from frame-counter gaps. A missing frame is a \emph{system} loss if no sensor at all reached the database in its 10-minute bin (network server, internet or the operators' logging pipeline), a \emph{gateway} loss if all of the sensor's serving gateways were silent, and a \emph{radio} loss otherwise. The classification is not an artefact of the 10-minute bin: with bins of 5 to 30~min the number of system outages ranges from \uvaBinOutMax{} to \uvaBinOutMin{} (short ones merge), the share of system down-time in outages of a day or more from \uvaBinDayMin{} to \uvaBinDayMax{}, and the system share of missing frames from \uvaBinSysMin{} to \uvaBinSysMax{}.

\emph{LoED.} The devices belong to other operators, so a frame counts as lost when no LoED gateway heard it, and most devices are simply far away. Only \emph{well-heard} devices are therefore used, those with at least 200 frames and a radio loss below 20\% (\loedWellHeard{} of \loedDevicesTwoHundred{}), as a proxy for a node's own link; a missing frame is a gateway loss if none of the device's serving gateways logged anything in that hour. Losses did not coincide across devices beyond what their density implies: a lost frame coincided with a loss of another device at a shared gateway within the same minute in \loedCoLoss{} of cases, against \loedCoLossBase{} (\loedCoLossLo--\loedCoLossHi) under a random time-shift baseline.

\emph{Frame counters.} Both LoRaWAN analyses count a frame once per device and frame counter, whichever gateways heard it and however often it was retransmitted. A counter decrease or a jump beyond a threshold (5{,}000 frames on UVA, 1{,}000 on LoED) is treated as a reset or rejoin, not as loss. On UVA, the frame counter decreased \uvaDecreases{} time in the whole record, no (sensor, counter) pair was received at two different times (no counter reuse), every sensor kept its 5-minute cadence in every month, and a threshold of 1{,}000 instead of 5{,}000 changes the count of missing frames by \uvaMissThouVsFive. On LoED, \loedResetDevices{} of the \loedWellHeardN{} well-heard devices show a reset, and their pooled loss is \loedLossHundred{} with a threshold of 100 and \loedLossThousand{} with 1{,}000.

\emph{Weija: the node's other streams.} The node's sensors share one logger and one uplink. A conductivity slot in which another stream delivered a genuine reading (not imputed, at most three decimals, not on an interpolated line) had a live node and a working uplink. Of the \weijaEcMissing{} missing conductivity slots, \weijaNodeSilentShare{} were node-silent and \weijaOtherStreamShare{} were conductivity-only gaps with the link up. In gaps of 6~h or more only \weijaLongAliveShare{} of the slots had another stream delivering. The node-level record (a slot is delivered if any stream delivered) loses \weijaNodeLoss{} of the slots in \weijaNodeGaps{} gaps, with \weijaNodeDayLost{} of the missing slots in gaps of a day or more, close to the conductivity record's \weijaCondDayLost{}.

\emph{InWaterSense: the second node.} Both nodes reached the server through one gateway with one GPRS link, housed in node~2's manhole~\cite{ahmedi2022dataset}. While node~2 was operating (May--August 2015, a window that holds \inwsWindowShare{} of node~1's missing slots), a node-1 gap during which node~2 kept delivering cannot be a gateway, GPRS or server outage. Node~2 was silent during \inwsNodeTwoSilent{} of node~1's missing slots. The share of node-1 missing slots with a shared cause, $\pi$, follows from $P(\text{node 2 silent}\mid\text{node-1 slot missing})=\pi+(1-\pi)\,b$, where $b$ is node~2's own rate of silence, measured within a day around each gap at slots where node~1 delivered. Gaps of 6~h or more were about \inwsSharedLong{} shared-cause (gateway, backhaul, server or a shared power failure) and \inwsSpecificLong{} node-specific; gaps shorter than 6~h were \inwsSpecificShort{} node-specific, as radio losses would be.

\section{Gap-length models}
\label{sec:s-models}
Besides the geometric model and the two-geometric mixture of the main paper, two heavy-tailed models are fitted to the gap lengths $k\ge1$:
\begin{align}
P_\mathrm{ln}(k)&=\frac{\Phi\!\left(\frac{\ln(k+\frac12)-\mu}{\sigma}\right)-\Phi\!\left(\frac{\ln(k-\frac12)-\mu}{\sigma}\right)}{1-\Phi\!\left(\frac{\ln\frac12-\mu}{\sigma}\right)}, \label{eq:s-ln}\\
P_\mathrm{pl}(k)&=\frac{k^{-\alpha}}{\zeta(\alpha)}, \label{eq:s-pl}
\end{align}
a lognormal discretised to integer lengths (\ref{eq:s-ln}), with $\Phi$ the standard normal distribution function, and a discrete power law (\ref{eq:s-pl}), with $\zeta$ the Riemann zeta function. All four are fitted by maximum likelihood and compared by Akaike's information criterion~\cite{akaike1974aic}.

\emph{What is tested.} The geometric distribution of gap lengths follows whenever the loss process has a single loss state with a constant chance of ending in each slot: independent (Bernoulli) losses, and the two-state Gilbert--Elliott channel, whose losses are correlated in time but whose loss periods still have geometric lengths. It is this single-loss-state assumption on gap lengths that is tested. The two-geometric mixture is a Markov chain with one up state and two loss states (``short loss'' and ``long gap'').

\emph{Absolute fit.} Because AIC only ranks models, each is also tested by a discrete Kolmogorov--Smirnov statistic $D=\max_k|\hat F(k)-F(k)|$ with a parametric-bootstrap $p$-value: 200 samples from the fitted model, each refitted. The geometric model is rejected on every record ($p\le$\,\gofPmin, the smallest value the 200 samples can give). With \gofInwsN{} gaps on InWaterSense and \gofUvaN{} runs on UVA, every tested model is rejected, by very different margins: the geometric model's $D$ is \gofGeomDInws{} on InWaterSense and \gofGeomDUva{} on UVA, against \gofBestDInws{} and \gofBestDUva{} for the best of the others. On Weija's few gaps, the lognormal and power-law models are not rejected ($p=$\,\gofWeijaLnP{} and \gofWeijaPlP). Without the longest gap of each record, every record keeps its best model. A gap is bounded by delivered slots on both sides, so every gap length is fully observed; only the first and last up-runs are censored, and they are not modelled.

\emph{Dependence between gaps.} On InWaterSense a long gap tends to be followed by another long gap (Spearman $\rho=$\,\inwsRhoNextOut) and by a short up-run ($\rho=$\,\inwsRhoNextUp; both \inwsRhoP), so gaps cluster into bad periods that a renewal model, drawing each gap independently, cannot reproduce. On Weija, with \weijaCondOutages{} gaps, neither correlation is distinguishable from zero ($\rho=$\,\weijaRhoNextOut, $p=$\,\weijaPNextOut; $\rho=$\,\weijaRhoNextUp, $p=$\,\weijaPNextUp). For these reasons all simulations replay the real records rather than a fitted model.

\section{Energy model and assumptions}
\label{sec:s-energy}
\emph{NB-IoT.} The measurements~\cite{michelinakis2021dissecting} send UDP datagrams of 12--512~B in normal coverage to a server that echoes them, without an application protocol such as CoAP. An unconfirmed report costs \eNbUnconf~J (connection released after the uplink) and a confirmed one \eNbConf~J (released after the server's echo, which is the acknowledgement, end to end). A bulk frame of up to 512~B, the module's largest datagram, costs up to 60\% more than a 20-byte report. The energy of a cell search when there is no network has no published measurement that could be found; it is modelled as the receive power of \pRxNb~W~\cite{sorensen2022nbiot} times a search time $T_\mathrm{search}$, swept over $\{2.2, 10, 60\}$~s (default 10~s, \eNbSearch~J per failed attempt).

\emph{LoRaWAN.} Energies come from the state-by-state current profile of~\cite{casals2017lorawan}, with airtime computed from the payload. At SF7 the acknowledgement arrives in the first receive window, so a confirmed success (\eSfSevenConf~J for a 20-byte report) is cheaper than an unconfirmed report (\eSfSevenUnconf~J), which keeps both windows open; at SF12 it arrives in the second window and a confirmed success is dearer (\eSfTwelveConf{} vs \eSfTwelveUnconf~J). A failed confirmed attempt adds RETRANSMIT\_TIMEOUT, 1--3~s~\cite{lora2021rp002}. Frames carry at most the regional maximum payload~\cite{lora2021rp002}, and the 1\% duty cycle forbids a new frame before 100 times the previous frame's airtime has passed.

\emph{Frames, host and floor.} A reading in a bulk frame takes 12~B (four parameters, a time offset and a sequence number) and a frame 4~B of overhead, so a frame holds \rpfNb{} readings on NB-IoT, \rpfSfSeven{} at SF7 and \rpfSfTwelve{} at SF12. The host is awake for every attempt, at $P_\mathrm{on}=\hostOnMW$~mW~\cite{libelium2020plugsense}. The lifetime $Y$ of the main paper uses the floor
\begin{equation}
E_\mathrm{floor}=86{,}400\,(P_\mathrm{sleep}^\mathrm{radio}+P_\mathrm{sleep}^\mathrm{host})+\frac{86{,}400}{T_\mathrm{s}}\,P_\mathrm{on}\,T_\mathrm{sense},\label{eq:s-floor}
\end{equation}
for sampling interval $T_\mathrm{s}$. In (\ref{eq:s-floor}), the radio sleeps at \nbPsmUW~$\mu$W (NB-IoT in power-saving mode~\cite{michelinakis2021dissecting}) or \loraSleepUW~$\mu$W (LoRa radio~\cite{casals2017lorawan}), the host sleeps at \hostSleepMW~mW~\cite{libelium2020plugsense}, and the host is assumed to be awake for $T_\mathrm{sense}=\tSenseS$~s per sample. The battery holds 6{,}600~mAh at \vNode~V ($B=\batteryKJ$~kJ).

\emph{Retry timing.} A record shows only that the reading at a slot's first instant was lost. Retries happen at their real times, after the receive windows, the retransmission time-out and the duty-cycle wait. In the \emph{full-slot} model the link stays down for the whole slot. In the \emph{restoration} model the link returns at a uniformly random time between the last missing and the next delivered instant of each gap, so only a gap's last slot can gain from a retry. In the \emph{burst} model, in addition, gaps of up to 1~h are radio-loss bursts in which every retry fails with the probability of a loss following a loss on LoED (\loedCondLoss), rather than surely.

Table~\ref{tab:s-assumptions} lists every model input.

\begin{table}[t]
\centering
\caption{Model inputs: measured (M) in the cited work, inferred (I) from the records, or assumed (A); where assumed, how it is tested.}
\label{tab:s-assumptions}
\scriptsize
\setlength{\tabcolsep}{2pt}
\begin{tabular}{@{}p{2.4cm}p{3.6cm}cp{1.9cm}@{}}
\toprule
Input & Value & & Tested by \\
\midrule
Link state & from the delivery records & I & gap cause, timing models \\
Retry timing in a gap & full slot / restoration / burst & A & all three \\
Reading generated in a gap & yes & A & node-off share 0--50\% \\
Radio loss & 0 or 0.096 per attempt (LoED) & I & both \\
Cadence & 611 s, 32.5 min, 5--60 min (USGS) & I & -- \\
Payload & 20 B report; 12 B per reading + 4 B & A & -- \\
ACK layer & LoRaWAN: network server; NB-IoT: application & M & application ACK variant \\
Transmissions & 8 (LoRaWAN 1.0.2), NbTrans 3 or 8 (1.0.4); NB-IoT 3 & M & both versions \\
LoRaWAN energies & state currents~\cite{casals2017lorawan} & M & $\times$0.5--2 \\
NB-IoT energies & per report~\cite{michelinakis2021dissecting} & M & $\times$0.5--2 \\
NB-IoT search & 0.222 W $\times$ 2.2--60 s & A & sweep \\
Host & 69 mW awake, 0.12 mW asleep~\cite{libelium2020plugsense} & M & $\times$0.5--2 \\
Sensing & 5 s awake per sample & A & ratios unaffected \\
Storage & unlimited, free & A & RAM 2--64 KB, flash \\
Backlog drain & at once & A & duty-cycle limit \\
Back-off cap $K$ & 6 h & A & 1--24 h \\
Battery & 6{,}600 mAh, ideal & A & ratios only \\
\bottomrule
\end{tabular}
\end{table}

\section{Loss handling: further results}
\label{sec:s-handling}
\emph{ARQ.} Under the full-slot model ARQ costs \arqEightMin--\arqEightMax{} the energy of unconfirmed reporting on LoRaWAN with 8 transmissions and \arqThreeMin--\arqThreeMax{} with 3; on InWaterSense at SF7 this is \inwSevenArqE~J/day instead of \inwSevenUnconfE~J/day (modelled lifetime \inwSevenArqLife{} instead of \inwSevenUnconfLife{} years). With independent radio losses added, ARQ recovers them, raising the delivered share by \arqGainMin{} to \arqGainMax{} percentage points. With eight transmissions at a fixed data rate (LoRaWAN 1.0.4, NbTrans~=~8) ARQ costs \arqVOneZeroFour{} the energy of unconfirmed reporting on InWaterSense at SF7, against \arqVOneZeroTwo{} with the stepping of LoRaWAN 1.0.2; at SF12 there is no lower rate to step to. NB-IoT ARQ recovers nothing in any timing model, because the modem reports that there is no network and the node stops retrying.

\emph{NB-IoT back-off.} Searching for a cell in every slot of a gap costs \backoffMin--\backoffMax{} the energy of backing off (unconfirmed reporting without back-off against unconfirmed reporting with it), the ratio growing with the share of time in gaps and with $T_\mathrm{search}$; with back-off a node in a long gap searches at most once every 6~h.

\emph{Buffering.} At SF7 buffering's energy relative to the fair baseline ranges from \bufSevenMin\% to \bufSevenMax\%, because a confirmed success costs less than an unconfirmed report and the node stops transmitting into gaps; at SF12 it is always dearer, by \bufTwelveMin\% to \bufTwelveMax\%. On InWaterSense at SF7 it uses \inwSevenBufE~J/day (\inwSevenBufLife~years). Under the other timing models its energy stays between \arqBufRatioMin{} and \arqBufRatioMax{} times that of unconfirmed reporting on LoRaWAN.

\emph{Node-off gaps.} When up to half of the slots in long gaps are reassigned to a node that was off, buffering still delivers every reading generated, but only \offExpectedInws{} of the expected readings on InWaterSense and \offExpectedWeija{} on Weija; its energy relative to the fair baseline moves from \offBufFairZero{} to \offBufFairHalf{} (SF7, InWaterSense). Weija has few long gaps, so its realised node-off shares are coarse (\offRealisedWeija{} for a nominal 0.1).

\emph{Storage and drain.} The longest backlog is \maxBufferKB~KB. A buffer in 2~KB of RAM loses \ramLossTwo{} of InWaterSense's readings, 8~KB loses \ramLossEight, 64~KB none. Writing every reading to serial flash~\cite{winbond2018w25q128jv}, with the host awake, costs at most \flashJDayMax~J/day in the datasheet's worst case, against at least \radioJDayMin~J/day of radio energy. If the backlog may use only what the duty cycle leaves after one live frame per slot, every reading is still delivered; at SF7 the longest backlog clears in \drainSevenMaxH~h; at SF12 the median backlog clears in \drainTwelveMedH~h, but in InWaterSense's worst period a backlog persists for \drainTwelveMaxD~days and readings generated while it drains arrive up to \drainTwelveLiveP~days late (95th percentile).

\emph{Back-off cap.} On InWaterSense at SF7, probing in every slot costs \kEverySlotSeven{} the energy of $K=6$~h (\kEverySlotNb{} on NB-IoT). Raising $K$ from 6 to 24~h saves a further \kSixToDaySave, while the 90th percentile of the delay between the link's return and the arrival of the gap's last reading grows from \kPninetySix~h to \kPninetyDay~h.

\emph{Many nodes.} The discrete-event model has one EU868 gateway and $N$ nodes that buffered through a common outage. Uplinks on the same channel and spreading factor collide unless one is 6~dB stronger; the gateway demodulates at most eight at once and cannot receive while it transmits; it acknowledges in the first receive window if its 1\% sub-band allows, else in the second (869.525~MHz, SF12, 10\% sub-band), else not at all; nodes back off as in the main paper, drain with up to eight transmissions per frame, and respect their own duty cycle. The collision rule reproduces the pure-ALOHA success probability $e^{-2G}$ in a test. Table~\ref{tab:s-netsim} gives the results.

\begin{table}[t]
\centering
\caption{$N$ nodes recovering from a common one-day outage behind one gateway (probe jitter up to 1~h; mean of 3 runs; 72~h horizon). Cleared: nodes whose backlog was delivered; median and last clearance (h after the outage); PDR: uplinks received and acknowledged; collided; received but unacknowledged; node energy over the horizon relative to a node recovering alone.}
\label{tab:s-netsim}
\scriptsize
\setlength{\tabcolsep}{2.4pt}
\begin{tabular}{@{}lrrrrrrrr@{}}
\toprule
 & $N$ & Cleared & Med.\ (h) & Last (h) & PDR & Coll. & No ACK & Energy \\
\midrule
\tabin{tables/tab_netsim.tex}
\bottomrule
\end{tabular}
\end{table}

\emph{Energy uncertainty.} The ratios were recomputed \mcDraws{} times with every measured component multiplied by an independent factor between 1/2 and 2 (log-uniform): unconfirmed and confirmed LoRaWAN reports, the acknowledgement time-out, the NB-IoT reports, bulk-size slope and search power, and the host's awake power, with the NB-IoT search time drawn between 2.2 and 60~s. On LoRaWAN, ARQ costs more than unconfirmed reporting in \mcArqDearer{} of the draws, at least \mcArqLow{} in 97.5\% of them; on NB-IoT it costs about the same (median \mcNbArqMed). Hourly batching is cheaper than the fair baseline in at least \mcBatchCheapOther{} of the draws at SF7 and on NB-IoT, but only in \mcBatchCheapTwelveMin--\mcBatchCheapTwelveMax{} at SF12, where it can cost up to \mcBatchTwelveHi{} as much. Buffering's energy relative to the fair baseline has a median of \mcBufMedMin--\mcBufMedMax{} but ranges from \mcBufLo{} to \mcBufHi{}: its sign is set by the ratio of confirmed to unconfirmed energy.

\section{Reconstruction}
\label{sec:s-reconstruct}
All series are placed on a common 1-h grid. For every ordered pair, parameter and gap length (\chainGaps~h), gaps are simulated at the target (a median of \chainOutages{} per pair, parameter and length), drawn from the last 80\% of the year. A gap covers the slots after the last observed slot $t_0$ and before the next observed slot $t_1$. Let $y$ be the neighbour's series, lagged by $\ell$ hours, and $b$ a scaling slope; $\ell\in[0,48]$~h maximises the correlation between 6-h changes of target and lagged neighbour in the first 20\% of the year, and $b$ is the least-squares slope of those changes (pairs whose best correlation is negative are dropped). Real-time estimators use only data up to $t_0$, and the neighbour during the gap:
\begin{align}
\hat{x}_\mathrm{pers}(t)&=x(t_0), \label{eq:s-pers}\\
\hat{x}_\mathrm{diel}(t)&=x(t_0)+x(t-24\,\mathrm{h})-x(t_0-24\,\mathrm{h}), \label{eq:s-diel}\\
\hat{x}_\mathrm{nb}(t)&=x(t_0)+b\,[y(t)-y(t_0)], \label{eq:s-nb}\\
\hat{x}_\mathrm{dn}(t)&=\hat{x}_\mathrm{diel}(t)+b\,[y(t)-y(t_0)] \nonumber\\
&\quad-b\,[y(t-24\,\mathrm{h})-y(t_0-24\,\mathrm{h})]. \label{eq:s-dn}
\end{align}
Persistence (\ref{eq:s-pers}) holds the last value; own diel (\ref{eq:s-diel}) is the server's estimator of the main paper; neighbour change (\ref{eq:s-nb}) adds the neighbour's scaled change since the gap began; and diel + neighbour (\ref{eq:s-dn}) adds only the part of that change which the neighbour's own diel pattern does not explain. Bridging estimators also know $x(t_1)$:
\begin{align}
\hat{x}_\mathrm{int}(t)&=x(t_0)+\frac{t-t_0}{t_1-t_0}\,[x(t_1)-x(t_0)], \label{eq:s-int}\\
\hat{x}_\mathrm{in}(t)&=\hat{x}_\mathrm{int}(t)+b\,[y(t)-\bar{y}(t)], \label{eq:s-in}
\end{align}
where $\bar{y}$ is the straight line from $y(t_0)$ to $y(t_1)$, so that (\ref{eq:s-in}) adds the neighbour's departure from its own line to the interpolation (\ref{eq:s-int}). Relative gains in the main paper exclude the \chainGainExcluded{} pH pairs whose target (site \chainExclSite) has a median error of exactly zero.

\emph{Spatial structure.} The neighbours lie \pairKmMin{} to \pairKmMax~km from their targets (median \pairKmMed~km); the order along the river is taken from the drainage area in the USGS site record, and for three sites without one from their position on the river. If the neighbour carried water to the target, the selected lag would be positive for upstream neighbours and zero for downstream ones; instead, for DO a downstream neighbour gets a positive lag in \lagPosDown{} of pairs and an upstream one in \lagPosUp.

\emph{Deep imputers.} SAITS~\cite{du2023saits} and StemGNN~\cite{cao2020stemgnn} (PyPOTS implementations) see all four parameters of every site of a chain in 192-h windows, normalised per window and series from the observed values. They are trained only on the first half of 2023, with the library's random point masks and, in half of the windows, a whole site masked for 1--72~h; the checkpoint with the best error on whole-site blocks in the last 15\% of the first half is kept. StemGNN uses one spectral block, one stack and width 16 (7.5~million parameters; its layers grow with the square of the window length). A first configuration, normalised with half-year means and spreads, under-fitted by an order of magnitude even on 1-h gaps and was discarded before any comparison. The models are tested on \deepGaps{} gap--parameter cases from the second half of the year, in which a whole site is missing for 1, 6, 24 or 72~h, and the paper's estimators are scored on the same cases, with the neighbour chosen and fitted on the first half. GRIN and SPIN need a graph library that the available machine could not hold, and ImputeFormer did not finish training in the available time.

\section{Design evaluation: protocol and further results}
\label{sec:s-design}

\begin{figure*}[t]
\centering
\includegraphics[width=\textwidth]{fig_mechanism}
\caption{When each reading reaches the server, in an illustrative scenario run through the simulator: LoRaWAN SF7, 15-min sampling, the link down from 6 to 10~h, and node alarms at 2.25~h (link up) and 7.5~h (link down). Points on the dotted diagonal arrive at once; horizontal runs are readings delivered together. Event-driven batching follows hourly batching except that the alarm at 2.25~h sends its batch early.}
\label{fig:s-mechanism}
\end{figure*}

\emph{Mechanism.} Fig.~\ref{fig:s-mechanism} shows the design's behaviours in a short scenario. Before the outage, the alarm makes the event-driven node upload at once, where plain batching would wait for the end of the hour; the alarm during the outage cannot get through; and after the outage, when the backlog arrives depends on where each strategy is in its back-off.

\emph{Development phase and deviations.} Hypotheses H1--H4 and their decision rules were written down before the runs for event-driven batching; after H4 failed, H5--H7 were written for the bounded-error trigger before it was implemented, and H8--H10 before they were run. Every deviation was recorded when made: UVA replayed as system outages plus radio losses (D1); one site excluded because its eligibility had been computed over its record span rather than the year (D2); seeds derived from site identifiers (D3); two implementation defects of the packing order, which made H8 fail as written (D4); and a floating-point tie at the tolerance (D5, found during the revision): pH and T are recorded in 0.1 steps, and an error of exactly one tolerance was counted in or out by rounding noise, in the scores and in the trigger. In the first development run, H4 also failed for pH (\hFourPassVOne{} combinations passing); those failures were D5, not batching.

\emph{Sites and links.} Sampling steps are 5 to 60~min; the first 20\% of each record trains the detector and all metrics use the remaining 80\%. The records are mapped onto each site's time grid in physical time, wrapped circularly with a random start offset per site (seeded by the site identifier, with a different seed in 2024). UVA is replayed as its system outages plus independent radio losses at the rate of a representative sensor (\uvaSensorFourRadio{} per attempt).

\emph{Injected events.} For parameter $j$, let $a_j(t)$ be the reading minus its centred 24-h running median, and $r_j(t)$ that anomaly minus its median for the same calendar month and hour of day. The event scale is
\begin{equation}
\sigma_j=\max\left(1.4826\,\operatorname{MAD}\!\left[r_j\right],\ \rho_j\right),
\label{eq:s-scale}
\end{equation}
floored at the recording resolution $\rho_j$, the most common positive difference between sorted distinct values. An event of amplitude $\kappa\sigma_j$, with $\sigma_j$ from (\ref{eq:s-scale}), $\kappa\in\{5,8\}$ and a random sign, lasts 180~min: a step adds $\kappa\sigma_j$ throughout, a ramp rises linearly from 0 to $\kappa\sigma_j$; values are rounded to $\rho_j$. Events touch one parameter each, in turn, lie in the test period, at least a day after the start of a data segment and wholly inside it, and at least three days apart (at most 60 per site and setting; \designEventsConf{} in 2024). Their median amplitudes are mild in physical units (5$\sigma$/8$\sigma$: T \ampFiveT/\ampEightT~$^\circ$C, DO \ampFiveDO/\ampEightDO~\%sat, EC \ampFiveEC/\ampEightEC~$\mu$S/cm, pH \ampFivePh/\ampEightPh), and real 3-hour changes of at least 5$\sigma$ occur at \realFiveT{} of the sites for T, \realFiveDO{} for DO, \realFiveEC{} for EC and \realFivePh{} for pH (8$\sigma$: \realEightT, \realEightDO, \realEightEC, \realEightPh). Of the 2024 events, \plausibleConf{} stay within T $-0.5$--40~$^\circ$C, positive EC, pH 4--11 and DO 0--250~\%sat (\plausibleDev{} in 2023), and the latency verdicts also hold on the plausible events alone.

\emph{Detector.} It compares each reading with the reading $m=120$~min earlier, corrected by a learned hour-of-day profile $p(\cdot)$ of the departure from the running 24-h median, and runs a two-sided CUSUM~\cite{page1954continuous} on the clipped, standardised residual:
\begin{align}
z_t&=\operatorname{clip}\!\left(\frac{x_t-x_{t-m}-p(h_t)+p(h_{t-m})}{s},\,-c,\,c\right),\label{eq:s-z}\\
S^{+}_t&=\max\left(0,\,S^{+}_{t-1}+z_t-k\right),\quad S^{-}_t=\max\left(0,\,S^{-}_{t-1}-z_t-k\right),\label{eq:s-cusum}
\end{align}
where $h_t$ is the hour of day and $s$ a robust residual scale floored at $\rho_j$. An alarm is raised when either sum in (\ref{eq:s-cusum}), driven by (\ref{eq:s-z}), exceeds $H$, with $k=3$, $H=8$ and $c=6$; the detector then stays in an event state until $|z_t|<2$ for 60~min. Profile and scale are initialised on the training period and updated by exponential smoothing (memory about 7~days) outside the event state. These parameters were fixed before any evaluation. In 2024 the detector caught a median (over sites) of \detectStepFive{} of 5$\sigma$ steps, \detectStepEight{} of 8$\sigma$ steps, \detectRampFive{} of 5$\sigma$ ramps and \detectRampEight{} of 8$\sigma$ ramps, with a median of \faMedian{} false alarms per parameter per 30 days (mean \faMean).

\emph{All injected events.} Counting every injected event, detected or not, the server knows of \knownOneEventNoneConf{} within an hour with event-driven batching on the no-outage link at SF7, the same as with per-reading reporting, against \knownOneBatchNoneConf{} with hourly batching; within a day it knows of \knownDayEventNoneConf{} with any of them. On the InWaterSense link the shares within an hour are \knownOneBoundedInwsConf{} (bounded), \knownOneArqInwsConf{} (ARQ + buffer), \knownOneUnconfInwsConf{} (unconfirmed) and \knownOneBatchInwsConf{} (batch), and within a day \knownDayBoundedInwsConf, \knownDayArqInwsConf, \knownDayUnconfInwsConf{} and \knownDayBatchInwsConf.

\emph{Picture.} Table~\ref{tab:s-picture} gives the picture for every parameter. Under the two water-quality gap records, unconfirmed reporting shows a better picture than ARQ + buffer for every parameter and radio (DO on InWaterSense at SF7: \picInwsDoUnconf{} vs \picInwsDoArq): when the link returns, an unconfirmed node delivers its next reading at once, while a buffering node waits for its next back-off probe.

\begin{table}[t]
\centering
\caption{Held-out year (2024): share of scored slots within the L2 tolerance (median over sites; LoRaWAN SF7; series without injected events). ARQ: ARQ + buffer.}
\label{tab:s-picture}
\scriptsize
\setlength{\tabcolsep}{3pt}
\begin{tabular}{@{}llrrrrr@{}}
\toprule
Link & Param. & Unconf. & ARQ & Batch & Event & Bounded \\
\midrule
\tabin{tables/tab_picture.tex}
\bottomrule
\end{tabular}
\end{table}

\emph{Events that start during a gap.} Table~\ref{tab:s-onset} and Fig.~\ref{fig:s-onset} show the detected events whose onset fell in a gap. Unconfirmed reporting learns of almost none (at most \onsetUnconfKnownMax{} on any link and radio); the buffering strategies learn of at least \onsetBufKnownMin, after the gap.

\begin{table}[t]
\centering
\caption{Held-out year (2024): detected events whose onset fell in a gap (LoRaWAN SF7, pooled over sites): number of events, share the server ever learns of, and median latency of those it learns of (hours).}
\label{tab:s-onset}
\scriptsize
\setlength{\tabcolsep}{2.4pt}
\begin{tabular}{@{}lrrrrrrrr@{}}
\toprule
 & & \multicolumn{3}{c}{Known ever} & \multicolumn{4}{c}{Median latency (h)} \\
\cmidrule(lr){3-5}\cmidrule(l){6-9}
Link & $n$ & Unconf. & ARQ & Event & ARQ & Batch & Event & Bounded \\
\midrule
\tabin{tables/tab_onset.tex}
\bottomrule
\end{tabular}
\end{table}

\begin{figure*}[t]
\centering
\includegraphics[width=\textwidth]{fig_outage_onset}
\caption{Held-out year (2024): detected events whose onset fell in a gap, share known at the server by the time on the horizontal axis (LoRaWAN SF7, pooled over sites; markers at 1~day).}
\label{fig:s-onset}
\end{figure*}

\emph{Packing order (H8).} On SF12 a frame holds only \rpfSfTwelve{} readings, so an hourly batch at 15-min sampling needs two frames, and the duty cycle delays the second. With first-in-first-out packing the newest reading rode in the second frame, and in the first development run the picture on the no-outage link fell to \fifoMinSfTwelve{} at worst. As written, H8 failed at one site (\hEightPreSite, 5-min sampling, SF12; minimum \hEightPreMin). The cause was two implementation defects: the reading promoted to the first frame could be a placeholder for a sensor gap, and the node's copy of the server counted readings as known as soon as their upload began, although at 5-min sampling an SF12 upload spans several slots. With both fixed (D4), newest-first packing is within tolerance at every scored slot of the no-outage link, on every site, parameter and radio, in both years, at \nfRatioMin--\nfRatioMax{} times the energy of first-in-first-out packing.

\emph{Tolerance level (H9, H10).} Table~\ref{tab:s-sensitivity} gives the energies at the three tolerance levels.

\begin{table}[t]
\centering
\caption{Held-out year (2024): energy (median J/day over sites, LoRaWAN SF7, series without injected events) at three tolerance levels, the median over sites of each site's ARQ / Bounded energy ratio, and the lowest within-tolerance share over the four parameters (median site). ARQ: ARQ + buffer.}
\label{tab:s-sensitivity}
\scriptsize
\setlength{\tabcolsep}{2.5pt}
\begin{tabular}{@{}llrrrrr@{}}
\toprule
 & & \multicolumn{2}{c}{Energy (J/day)} & & \multicolumn{2}{c}{Within tolerance} \\
\cmidrule(lr){3-4}\cmidrule(l){6-7}
Level & Link & ARQ & Bounded & ARQ/Bounded & Bounded & ARQ \\
\midrule
\tabin{tables/tab_sensitivity.tex}
\bottomrule
\end{tabular}
\end{table}

\emph{Tolerances and sensor accuracy.} L2 lies at or inside the U.S. Geological Survey's rating of an ``excellent'' continuous record~\cite{wagner2006usgs}: over the sites used, the median L2 tolerance is \tolTwoT~$^\circ$C against \usgsExcT~$^\circ$C for T, \tolTwoDO{} against \usgsExcDO~\%sat for DO (0.3~mg/L or 5\% of the reading), \tolTwoEC{} against \usgsExcEC~$\mu$S/cm for EC (3\% of the reading; USGS rates conductance in percent, as the EC tolerance is defined) and \tolTwoPh{} against \usgsExcPh{} for pH. All three levels are small compared with the median daily range at these sites (T \dayRangeT~$^\circ$C, DO \dayRangeDO~\%sat, EC \dayRangeEC~$\mu$S/cm, pH \dayRangePh).

\bibliographystyle{IEEEtran}
\bibliography{refs}

\end{document}
